\subsection{The self-dual Higgs submodule associated with a nilpotent element}
We explain in this section that an element $y\in\mathfrak m^-$ defines a special graded subspace in $\jepPubBbb{E}$. This algebraic construction will be used in Section~4.3 to produce the leafwise Higgs subsheaves needed to prove the Milnor-Wood inequality.
\subsubsection{$K$-orbits in $\mathfrak m^-$: ranks of nilpotent elements}
We begin by describing the $K$-orbits in $\mathfrak m^-$, introducing the notion of rank, and then we choose nice representatives in the orbits to simplify our work. Let us recall the following result, which is proved e.g. in \cite{jep93HC56}, \cite[Ch.~VIII, \S7]{jep93Hel01}, or \cite{jep93Wol72}.
\begin{proposition}
The orbits of the complex group $K$ in $\mathfrak m^-$ are parametrized by integers $r\in\{0,\ldots,r_M\}$. A representative of each orbit is $y_{\alpha_1}+\cdots+y_{\alpha_r}$, where $(\alpha_1,\ldots,\alpha_{r_M})$ is the maximal dominant orthogonal sequence.
\end{proposition}
\begin{definition}
Let $y\in\mathfrak m^-$. The rank $r(y)$ of $y$ is the integer $r$ such that $y$ is in the $K$-orbit of $y_{\alpha_1}+\cdots+y_{\alpha_r}$.
\end{definition}
\begin{example}
In case $\mathfrak g$ has type $A$, $\mathfrak m^-$ identifies with a space of matrices, and the rank as defined above of an element in $\mathfrak m^-$ coincides with its rank as a matrix.
\end{example}
The following proposition is a characterization of the rank $r(y)$ of $y\in\mathfrak m^-$ in terms of its action on $\jepPubBbb{E}$. If $y\in\mathfrak g$, we will denote by $\jepPubBmi{y}\in\operatorname{End}(\jepPubBbb{E})$ the image of $y$ by the representation of $\mathfrak g$ on $\jepPubBbb{E}$.
\begin{proposition}
Let $y=y_{\alpha_1}+\cdots+y_{\alpha_{r(y)}}$ with $(\alpha_1,\ldots,\alpha_{r(y)})$ a dominant orthogonal sequence. We have $\jepPubBmi{y}^{r(y)+1}=0$ and $\jepPubBmi{y}^{r(y)}(\jepPubBbb{E}_0)=\jepPubBbb{E}_{\varpi-\alpha_1-\cdots-\alpha_{r(y)}}$, so in particular $\jepPubBmi{y}^{r(y)}(\jepPubBbb{E}_0)\neq\{0\}$. In other words, the rank of $y$ is the order of nilpotency of $\jepPubBmi{y}$.
\end{proposition}
\begin{proof}
We have $y_{\alpha_i}\in\mathfrak g_{-\alpha_i}$. For any pair $(\chi,\delta)$ with $\chi\in X(\jepPubBbb{E})$ and $\delta$ a long root, $\chi-2\delta$ cannot be a weight of $\jepPubBbb{E}$ by Lemma~2.9. We deduce from Proposition~2.19~(1) that $\jepPubBmi{y}_{\alpha_i}^2=0$. On the other hand, for any $i,j$ with $i\neq j$, the maps $\jepPubBmi{y}_{\alpha_i}$ and $\jepPubBmi{y}_{\alpha_j}$ commute since $\alpha_i+\alpha_j$ is not a root by Proposition~2.19~(1) also. Thus we get
\[
\jepPubBmi{y}^k=k!\sum\jepPubBmi{y}_{\alpha_{j_1}}\circ\cdots\circ\jepPubBmi{y}_{\alpha_{j_k}},
\]
where the sum is over the increasing sequences $1\leqslant j_1<\cdots<j_k\leqslant r(y)$.

In particular $\jepPubBmi{y}^{r(y)+1}=0$, and $\jepPubBmi{y}^{r(y)}=r(y)!\jepPubBmi{y}_{\alpha_1}\circ\cdots\circ\jepPubBmi{y}_{\alpha_{r(y)}}$. By Lemma~2.9, we have $\jepPubBmi{y}_{\alpha_1}\circ\cdots\circ\jepPubBmi{y}_{\alpha_{r(y)}}\cdot\jepPubBbb{E}_\varpi=\jepPubBbb{E}_{\varpi-\alpha_1-\cdots-\alpha_{r(y)}}$, so the proposition is proved.
\end{proof}
For the remaining of this section, we let $1\leqslant r\leqslant r_M$ be an integer and $y=y^r:=y_{\alpha_1}+\cdots+y_{\alpha_r}\in\mathfrak m^-$, where $(\alpha_1,\ldots,\alpha_r)$ is the dominant orthogonal sequence of length $r$ obtained by the algorithm of Definition~2.17. By Proposition~2.21, there is no loss of generality in considering this particular element $y$. We denote by $h_r$ the element $\alpha_1^\jepPubCoroot+\cdots+\alpha_r^\jepPubCoroot$ spelled out in Table~1. Recall that $\jepPubBmi{y}$ is $y$ acting on $\jepPubBbb{E}$ via the cominuscule representation of $\mathfrak g$, and that $r$ is the order of nilpotency of $\jepPubBmi{y}$.
\subsubsection{Definition of $\jepPubBbb{V}^r$ and the Higgs property}
It is a very well-known idea that the nilpotent element $\jepPubBmi{y}$ of $\operatorname{End}(\jepPubBbb{E})$ uniquely defines an increasing filtration $\jepPubBbb{W}^r_\jepPubBullet$ of $\jepPubBbb{E}$, called the \emph{Deligne weight filtration} of $\jepPubBmi{y}$, or Jacobson-Morozov filtration. This idea appears in \cite[\S1.6.13]{jep93Del80} in the context of Hodge theory and in \cite[\S\S3.2 \& 3.4]{jep93McG02} in the context of nilpotent orbits in Lie algebras.

The defining properties of the filtration $(\jepPubBbb{W}^r_k)_{k\in\mathbb Z}$ are that for all $k\in\mathbb Z$,
\[
\jepPubBmi{y}(\jepPubBbb{W}^r_k)\subset\jepPubBbb{W}^r_{k-2}
\]
and for all $k\geqslant1$,
\begin{equation}\tag{2.1}
\jepPubBmi{y}^k:\jepPubBbb{W}^r_k/\jepPubBbb{W}^r_{k-1}\longrightarrow\jepPubBbb{W}^r_{-k}/\jepPubBbb{W}^r_{-k-1}\quad\text{is an isomorphism.}
\end{equation}
The filtration $(\jepPubBbb{W}^r_k)_{k\in\mathbb Z}$ can be defined as follows (see e.g. \cite{jep93SZ85}):
\begin{equation}\tag{2.2}
\jepPubBbb{W}^r_k=\sum_{\substack{\ell\geqslant0\\k+\ell+1\geqslant0}}\operatorname{Ker}\jepPubBmi{y}^{k+\ell+1}\cap\operatorname{Im}\jepPubBmi{y}^{\ell}=\sum_{\substack{i-j=k\\i\geqslant0,\ j\geqslant0}}\operatorname{Ker}\jepPubBmi{y}^{i+1}\cap\operatorname{Im}\jepPubBmi{y}^{j}.
\end{equation}
Since $r$ is the rank of nilpotency of $\jepPubBmi{y}$ acting on $\jepPubBbb{E}$, we have $\jepPubBbb{W}^r_k=\{0\}$ for $k<-r$ and $\jepPubBbb{W}^r_k=\jepPubBbb{E}$ for $k\geqslant r$.

The elements $y$ and $h_r$ fit in an $\mathfrak{sl}_2$-triple $(x,h_r,y)$ for some $x\in\mathfrak m^+$, and the subspaces $\jepPubBbb{W}^r_k$ can be alternatively defined by means of this triple:
\begin{equation}\tag{2.3}
\jepPubBbb{W}^r_k=\jepPubOplus_{\chi\in X(\jepPubBbb{E}):\langle\chi,h_r\rangle\leqslant k}\jepPubBbb{E}_\chi.
\end{equation}
We combine the gradation $\jepPubBbb{E}=\jepPubOplus_{i=0}^{r_M}\jepPubBbb{E}_i$ and the weight filtration $(\jepPubBbb{W}^r_k)_{k\in\mathbb Z}$ to define the subspaces of $\jepPubBbb{E}$ that will be the center of our attention for the remaining part of Section~2.
\begin{notation}
Let $\jepPubBbb{V}^r_i=\jepPubBbb{E}_i\cap\jepPubBbb{W}^r_{r-2i}$ for $0\leqslant i\leqslant r_M$, and let $\jepPubBbb{V}^r=\jepPubOplus_i\jepPubBbb{V}^r_i\subset\jepPubBbb{E}$.
\end{notation}
Observe that $\jepPubBbb{V}^r_0=\jepPubBbb{E}_0$ and that $\jepPubBbb{V}^r_i\neq\{0\}$ only if $i\leqslant r$. The subspaces $\jepPubBbb{V}^r_i$ can also be described as sums of weight spaces $\jepPubBbb{E}_\chi$.
\begin{proposition}
For $\chi\in X(\jepPubBbb{E}_i)$, we have $\langle\chi,h_r\rangle\geqslant r-2i$. Hence, for all $i$ such that $0\leqslant i\leqslant r$,
\[
\jepPubBbb{V}^r_i=\jepPubOplus_{\substack{\chi\in X(\jepPubBbb{E}_i):\\\langle\chi,h_r\rangle=r-2i}}\jepPubBbb{E}_\chi=\jepPubOplus_{\substack{\chi\in X(\jepPubBbb{E}):\\\langle\chi,z\rangle=z_{\mathtt{max}}-2i,\\\langle\chi,h_r\rangle=r-2i}}\jepPubBbb{E}_\chi.
\]
\end{proposition}
\begin{proof}
Proposition~2.15 and the fact that, for any root $\alpha$, we have $\langle\alpha,h_r\rangle\leqslant2$ (see Proposition~2.20) imply the first assertion. It is plain from the description of the subspaces $\jepPubBbb{W}^r_k$ given in~(2.3) that the second assertion follows from the first.
\end{proof}
The following lemma is reminiscent of Equation~(2.1). We will prove a stronger statement in Proposition~2.40. More Hodge-theoretic considerations will be made in Section~2.6.
\begin{lemma}
The linear map $\jepPubBmi{y}^{r-2i}$ is an isomorphism between $\jepPubBbb{V}^r_i$ and $\jepPubBbb{V}^r_{r-i}$. In particular, $\dim\jepPubBbb{V}^r_{r-i}=\dim\jepPubBbb{V}^r_i$.
\end{lemma}
\begin{proof}
We know that
\[
\jepPubBbb{V}^r_i=\jepPubBbb{E}_i\cap\jepPubBbb{W}^r_{r-2i},\qquad\jepPubBbb{V}^r_{r-i}=\jepPubBbb{E}_{r-i}\cap\jepPubBbb{W}^r_{2i-r},
\]
and that for $i$ satisfying $0\leqslant i\leqslant r/2$, $\jepPubBmi{y}^{r-2i}$ is an isomorphism between $\jepPubBbb{W}^r_{r-2i}/\jepPubBbb{W}^r_{r-2i-1}$ and $\jepPubBbb{W}^r_{2i-r}/\jepPubBbb{W}^r_{2i-r-1}$. By Proposition~2.15, $\jepPubBmi{y}^{r-2i}$ maps $\jepPubBbb{E}_i$ to $\jepPubBbb{E}_{r-i}$. Since we have $\jepPubBbb{E}_k\cap\jepPubBbb{W}^r_{r-2k-1}=\{0\}$ for all $k$ by Proposition~2.26 and Equation~(2.3), we get that $\jepPubBmi{y}^{r-2i}$ is an isomorphism between $\jepPubBbb{V}^r_i$ and $\jepPubBbb{V}^r_{r-i}$.
\end{proof}
\begin{definition}
A subspace $\jepPubBbb{S}$ of $\jepPubBbb{E}$ is a \emph{$y$-Higgs subspace} if $y\cdot\jepPubBbb{S}\subset\jepPubBbb{S}$ and $\mathfrak m^+\cdot\jepPubBbb{S}\subset\jepPubBbb{S}$.
\end{definition}
The following algebraic fact is at the heart of the construction of leafwise Higgs subsheaves in Section~4.
\begin{proposition}
The subspace $\jepPubBbb{V}^r$ is a $y$-Higgs subspace of $\jepPubBbb{E}$. More precisely, the following inclusions hold:
\begin{itemize}
\item $y\cdot\jepPubBbb{V}^r_i\subset\jepPubBbb{V}^r_{i+1}$
\item $\mathfrak m^+\cdot\jepPubBbb{V}^r_i\subset\jepPubBbb{V}^r_{i-1}$
\end{itemize}
\end{proposition}
\begin{proof}
Using the description of $\jepPubBbb{V}^r_i$ given in Proposition~2.26, the first inclusion holds because $y\in\jepPubOplus_i\mathfrak g_{-\alpha_i}$ and each $\alpha_i$ satisfies $\langle-\alpha_i,z\rangle=\langle-\alpha_i,h_r\rangle=-2$. The second one also holds since for $\alpha$ a root of $\mathfrak m^+$, we have $\langle\alpha,z\rangle=2$ (and $\langle\alpha,h_r\rangle\leqslant2$ by Proposition~2.20).
\end{proof}
\subsubsection{The module structure of $\jepPubBbb{V}^r$}
We now show that the subspace $\jepPubBbb{V}^r$ of $\jepPubBbb{E}$ has the structure of an irreducible $L_r$-module, for $L_r$ a complex reductive group determined by $y$, whereas the $\jepPubBbb{V}^r_i$ are irreducible $H_r$-modules, where $H_r=L_r\cap K$.

We begin by defining the relevant subalgebras of $\mathfrak g$.
\begin{notation}
Let $\mathfrak q_r\subset\mathfrak g$ be the parabolic subalgebra defined by the coweight $z-h_r$:
\[
\mathfrak q_r=\mathfrak t\oplus\jepPubOplus_{\alpha:\langle\alpha,h_r\rangle\leqslant\langle\alpha,z\rangle}\mathfrak g_\alpha.
\]
A Levi factor of $\mathfrak q_r$ is
\[
\mathfrak l_r=\mathfrak t\oplus\jepPubOplus_{\alpha:\langle\alpha,h_r\rangle=\langle\alpha,z\rangle}\mathfrak g_\alpha.
\]
Let $\mathfrak l_r^\pm\subset\mathfrak l_r$ be the nilpotent subalgebras of $\mathfrak l_r$ defined by
\[
\mathfrak l_r^\pm=\jepPubOplus_{\alpha:\langle\alpha,h_r\rangle=\langle\alpha,z\rangle=\pm2}\mathfrak g_\alpha.
\]
The intersection $\mathfrak k\cap\mathfrak q_r$ is a parabolic subalgebra of $\mathfrak k$ and a Levi factor of it is
\[
\mathfrak h_r=\mathfrak t\oplus\jepPubOplus_{\alpha:\langle\alpha,h_r\rangle=\langle\alpha,z\rangle=0}\mathfrak g_\alpha.
\]
We denote by $Q_r$, $L_r$ and $H_r$, the connected subgroups of $G$ with Lie algebra $\mathfrak q_r$, $\mathfrak l_r$ and $\mathfrak h_r$ respectively. Such groups exist because $\mathfrak q_r$ is a parabolic subalgebra, and $\mathfrak h_r$ and $\mathfrak l_r$ are Levi subalgebras. We observe that $H_r$ is a Levi subgroup of $K\cap Q_r$: in fact, we have $K\cap Q_r=H_r\cdot R(K\cap Q_r)$, where $R(K\cap Q_r)$ denotes the unipotent radical of $K\cap Q_r$. Since Levi subgroups correspond to Levi subalgebras, $L_r$ is a Levi factor of $Q_r$.
\end{notation}
For the convenience of the reader, the conditions that define the different Lie subalgebras of $\mathfrak g$ we are considering are displayed in Table~2.
\begin{table}[H]\centering
\begin{tabular}{lcccccc}\hline
Lie algebra & $\mathfrak k$ & $\mathfrak m^\pm$ & $\mathfrak q_r$ & $\mathfrak l_r$ & $\mathfrak l_r^\pm$ & $\mathfrak h_r$\\
Condition & $z=0$ & $z=\pm2$ & $h_r\leqslant z$ & $h_r=z$ & $h_r=z=\pm2$ & $h_r=z=0$\\\hline
\end{tabular}
\caption{\small Lie subalgebras under consideration (We abbreviate the condition $\langle\alpha,z\rangle=0$, resp. $\langle\alpha,h_r\rangle=0$, on a root $\alpha$ as $z=0$, resp. $h_r=0$.)}
\end{table}
We have the following lemmas concerning the $\jepPubBbb{V}^r_i$’s.
\begin{lemma}
We have $\jepPubBbb{V}^r_{i+1}=\mathfrak l_r^-\cdot\jepPubBbb{V}^r_i$ and $\jepPubBbb{V}^r_{i-1}=\mathfrak l_r^+\cdot\jepPubBbb{V}^r_i$.
\end{lemma}
\begin{proof}
Let us denote by $X(\jepPubBbb{E}_i)$ the set of weights of $\jepPubBbb{E}_i$. We know by Proposition~2.15 that $\jepPubBbb{E}_{i+1}=\mathfrak m^-\cdot\jepPubBbb{E}_i$. Thus,
\[
\jepPubBbb{E}_{i+1}=\biggl(\jepPubOplus_{\alpha:\langle\alpha,z\rangle=-2}\mathfrak g_\alpha\biggr)\cdot\biggl(\jepPubOplus_{\chi\in X(\jepPubBbb{E}_i)}\jepPubBbb{E}_\chi\biggr)=\jepPubOplus_{\substack{\chi\in X(\jepPubBbb{E}_i)\\\alpha\in R(\mathfrak m^-)}}\jepPubBbb{E}_{\chi+\alpha}.
\]
Given $\chi\in X(\jepPubBbb{E}_i)$ and $\alpha\in R(\mathfrak m^-)$, we can make two observations:
\begin{itemize}
\item If $\langle\chi,h_r\rangle>r-2i$ then we have $\langle\chi+\alpha,h_r\rangle>r-2(i+1)$ and thus $\jepPubBbb{E}_{\chi+\alpha}\not\subset\jepPubBbb{V}^r_{i+1}$
\item If $\langle\alpha,h_r\rangle>-2$ then $\langle\chi+\alpha,h_r\rangle>r-2(i+1)$ and thus $\jepPubBbb{E}_{\chi+\alpha}\not\subset\jepPubBbb{V}^r_{i+1}$
\end{itemize}
The first equality of the lemma follows from these observations. The proof of the second equality is similar.
\end{proof}
Note that $\mathfrak l_r^-$ is an $\mathfrak h_r$-module. More precisely, we have
\begin{lemma}
The modules $\jepPubBbb{V}^r_i$ are irreducible $\mathfrak h_r$-modules. The Lie algebra $\mathfrak l_r^-$ is an irreducible $\mathfrak h_r$-module.
\end{lemma}
\begin{proof}
Denote by $\mathfrak k^+$, resp. $\mathfrak h_r^+$, the sum of the root spaces for positive roots in $\mathfrak k$, resp. $\mathfrak h_r$. Let $i$ be fixed and such that $\jepPubBbb{V}^r_i\neq\{0\}$. By Proposition~2.15, $\jepPubBbb{E}_i$ is an irreducible $\mathfrak k$-module. Let $\mu_i$ be the lowest weight of $\jepPubBbb{E}_i$. We have $\jepPubBbb{E}_{\mu_i}\subset\jepPubBbb{V}^r_i$. Since $\jepPubBbb{E}_i$ is irreducible, we have $\jepPubBbb{E}_i=U(\mathfrak k^+)\cdot\jepPubBbb{E}_{\mu_i}$ (here $U$ denotes the universal enveloping algebra). Now, arguing as in the proof of the Lemma~2.31, we see that this implies that $\jepPubBbb{V}^r_i=U(\mathfrak h_r^+)\cdot\jepPubBbb{E}_{\mu_i}$. Now, $\jepPubBbb{E}_{\mu_i}$ has dimension 1, hence $\jepPubBbb{V}^r_i$ is indecomposable. Since $\mathfrak h_r$ is reductive, this proves that $\jepPubBbb{V}^r_i$ is irreducible. Now, by Lemma~2.31 again, we have $\jepPubBbb{V}^r_1=\mathfrak l_r^-\cdot\jepPubBbb{V}^r_0\simeq\mathfrak l_r^-\cdot\jepPubBbb{E}_0$: thus $\mathfrak l_r^-$ is also irreducible.
\end{proof}
Concerning the submodule $\jepPubBbb{V}^r$, we have
\begin{proposition}
\begin{enumerate}[label=(\arabic*)]
\item The subspace $\jepPubBbb{V}^r\subset\jepPubBbb{E}$ is an irreducible $Q_r$-module
\item As a consequence, the nilpotent radical of $Q_r$ acts trivially on $\jepPubBbb{V}^r$
\item $\jepPubBbb{V}^r_i$ is a $K\cap Q_r$-module
\item The subgroup of elements in $G$ which preserve $\jepPubBbb{V}^r$ is exactly $Q_r$
\end{enumerate}
\end{proposition}
\begin{proof}
In characteristic $0$, given a representation of a connected group, a subspace of this representation is stable under this group if and only if it is stable for the induced action of the Lie algebra. Thus, to prove the proposition, it is enough to deal with the corresponding Lie algebras.

Let us prove that $\jepPubBbb{V}^r$ is stable under $\mathfrak q_r$. It is clearly stable under $\mathfrak t$. Let $\alpha$ be such that $\mathfrak g_\alpha\subset\mathfrak q_r$, i.e., $\langle\alpha,h_r\rangle\leqslant\langle\alpha,z\rangle$. Let $v\in\jepPubBbb{E}_\chi\subset\jepPubBbb{V}^r_i$, so that we have $\langle\chi,h_r\rangle=r-2i$. For $x\in\mathfrak g_\alpha$, we have $x\cdot v\in\jepPubBbb{E}_{\chi+\alpha}\subset\jepPubBbb{E}_{i-\langle\alpha,z\rangle/2}$. Since $\langle\alpha,h_r\rangle\leqslant\langle\alpha,z\rangle$, we have $\langle\chi+\alpha,h_r\rangle\leqslant r-2i+\langle\alpha,z\rangle$, thus either $x\cdot v=0$ or $\langle\chi+\alpha,h_r\rangle=r-2i+\langle\alpha,z\rangle$, by Proposition~2.26. In the second case, we get $x\cdot v\in\jepPubBbb{V}^r_{i-\langle\alpha,z\rangle/2}$. Combining Lemmas~2.31 and~2.32, we see that $\jepPubBbb{V}^r$ is an irreducible $\mathfrak l_r$-module, hence also an irreducible $\mathfrak q_r$-module.

It follows from Engel’s theorem that there exists a non-zero vector in $\jepPubBbb{V}^r$ which is annihilated by the nilpotent radical of $\mathfrak q_r$. Since this radical is an ideal and $\jepPubBbb{V}^r$ is an irreducible $\mathfrak q_r$-module, this radical acts trivially on $\jepPubBbb{V}^r$.

The fact that $\jepPubBbb{V}^r_i$ is a $K\cap Q_r$-module follows because $\jepPubBbb{V}^r_i=\jepPubBbb{E}_i\cap\jepPubBbb{V}^r$, $\jepPubBbb{E}_i$ is $K$-stable, and $\jepPubBbb{V}^r$ is $Q_r$-stable.

Finally, to prove that the stabilizer of $\jepPubBbb{V}^r$ is exactly $Q_r$, let us denote by $\mathfrak{stab}(\jepPubBbb{V}^r)\subset\mathfrak g$ the Lie subalgebra preserving the subspace $\jepPubBbb{V}^r$ in $\jepPubBbb{E}$. We know by~(1) that $\mathfrak{stab}(\jepPubBbb{V}^r)\supset\mathfrak q_r$. On the other hand, let $\alpha$ be a root such that $\mathfrak g_\alpha\cdot\jepPubBbb{V}^r\subset\jepPubBbb{V}^r$. Let us assume as a first case that $\langle\alpha,z\rangle=-2$, i.e., $\mathfrak g_\alpha\subset\mathfrak m^-$. Since, by Proposition~2.15(d), the action of $\mathfrak g$ on $\jepPubBbb{E}$ induces an isomorphism $\jepPubBbb{E}_1\simeq\jepPubBbb{E}_\varpi\otimes\mathfrak m^-$, we have $\mathfrak g_\alpha\cdot\jepPubBbb{V}^r_0=\mathfrak g_\alpha\cdot\jepPubBbb{E}_\varpi=\jepPubBbb{E}_{\varpi+\alpha}$. Then $\jepPubBbb{E}_{\varpi+\alpha}\subset\jepPubBbb{V}^r\cap\jepPubBbb{E}_1=\jepPubBbb{V}^r_1$, so $\langle\alpha,h_r\rangle=-2$, and $\mathfrak g_a\subset\mathfrak q_r$. Assume now that $\langle\alpha,z\rangle=0$, i.e., $\mathfrak g_\alpha\subset\mathfrak k$, and by contradiction that $\langle\alpha,h_r\rangle>0$. Proposition~2.20 then implies that $\langle\alpha,h_r\rangle=1$. For any integer $1\leqslant i\leqslant r$, we cannot have $\langle\alpha,\alpha_i^\jepPubCoroot\rangle=2$ because this would imply $\alpha=\alpha_i$ and $\langle\alpha,h_r\rangle=2$. Let $i$ be such that $\langle\alpha,\alpha_i^\jepPubCoroot\rangle>0$: then $\langle\alpha,\alpha_i^\jepPubCoroot\rangle=1$. Therefore, $\alpha-\alpha_i$ is a root. Since $\mathfrak g_\alpha\subset\mathfrak k$, $\mathfrak g_\alpha\cdot\jepPubBbb{E}_i\subset\jepPubBbb{E}_i$ and hence $\mathfrak g_\alpha\cdot\jepPubBbb{V}^r_i\subset\jepPubBbb{V}^r_i$. Furthermore $\mathfrak g_\alpha(\jepPubBbb{V}^r_0)=0$ since $\jepPubBbb{V}^r_0=\jepPubBbb{E}_\varpi$. It follows that
\[
\mathfrak g_{\alpha-\alpha_i}\cdot\jepPubBbb{E}_\varpi=[\mathfrak g_\alpha,\mathfrak g_{-\alpha_i}]\cdot\jepPubBbb{E}_\varpi=\mathfrak g_\alpha\cdot(\mathfrak g_{-\alpha_i}\cdot\jepPubBbb{E}_\varpi).
\]
Again by Proposition~2.15(d), $\mathfrak g_{\alpha-\alpha_i}\cdot\jepPubBbb{E}_\varpi=\jepPubBbb{E}_{\varpi+\alpha-\alpha_i}$ and $\mathfrak g_{-\alpha_i}\cdot\jepPubBbb{E}_\varpi=\jepPubBbb{E}_{\varpi-\alpha_i}$. But $\jepPubBbb{E}_{\varpi-\alpha_i}\subset\jepPubBbb{V}^r_1$ whereas $\jepPubBbb{E}_{\varpi+\alpha-\alpha_i}\not\subset\jepPubBbb{V}^r_1$ since we assumed $\langle\alpha,h_r\rangle>0$. This contradicts $\mathfrak g_\alpha\cdot\jepPubBbb{V}^r\subset\jepPubBbb{V}^r$.

Let now $\operatorname{Stab}(\jepPubBbb{V}^r)\subset G$ be the subgroup stabilizing $\jepPubBbb{V}^r$. We have $Q_r\subset\operatorname{Stab}(\jepPubBbb{V}^r)$, thus $\operatorname{Stab}(\jepPubBbb{V}^r)$ is parabolic and therefore connected \cite[Cor.~23.1.B]{jep93Hum75}. It follows that $Q=\operatorname{Stab}(\jepPubBbb{V}^r)$.
\end{proof}

